\clearpage
\section*{章末問題}
\begin{enumerate}[label=\textbf{\thechapter.\arabic*},leftmargin=*]
 \item\label{問:導入速度} $s(t)=t^2$ とする.時刻 $3$ から $3+h$ までの平均の速度を求めよ.\,$h=0.2,-0.2$ の場合の値を求め,時刻 $3$ の瞬間の速度を予想せよ.
 \item\label{問:導入直線面積} $y=2x+1$ と $x$ 軸,直線 $x=0,x=2$ に囲まれた台形の面積を求めよ.また,区間 $[0,2]$ を $n$ 等分して左端・右端の高さから作る長方形和 $L_n,U_n$ を求め,それぞれが同じ面積に近づくことを確かめよ.
 \item\label{問:導入中点} $y=x^2$ の下の $0\le x\le1$ の面積を,\,$4$ 等分した各小区間の中央の高さで近似せよ.得られた値を $L_4,U_4$ と比較せよ.
 \item\label{問:導入三次割線} $y=x^3$ のグラフ上の2点 $(1,1)$,\,$(1+h,(1+h)^3)$ を結ぶ直線の傾きを求めよ.\,$h\ne0$ の条件を使う箇所を明記し,\,$h$ を $0$ に近づけたときの傾きを予想せよ.
 \item\label{問:導入拡大面積} $y=x^2$ の下の $0\le x\le2$ の面積を $S$ とする.\,$4$ 等分した左端と右端の長方形和を求めて $S$ を挟め.一般の $n$ 等分では,\,$0\le x\le1$ の場合の長方形和の何倍になるか.そこから $S$ を予想せよ.
 \item\label{問:導入符号} 直線上を動く物体の速度が,時刻 $0$ から $4$ までは毎秒 $3$ メートル,時刻 $4$ から $7$ までは毎秒 $-2$ メートルである.位置の変化と道のりをそれぞれ求めよ.2つの答えが異なる理由を,速度の符号を使って説明せよ.
 \item\label{問:導入二次速度} 時刻$t\ge0$での速度が$v(t)=t^2$である運動について,時刻$0$から$T$までを$n$等分した左端と右端の和を求めよ.そこから位置の変化を求めよ.ただし$T>0$とする.
 \item\label{問:導入三角形和} 取り尽くし法の本文で定めた$P_N$を求め,極限値$\frac16$との差を表せ.その差を$\frac1{1000}$未満にする最小の$0$以上の整数$N$を求めよ.
\end{enumerate}

\smallskip\noindent 解答は\pageref{解答:第1章}ページ.

